\secciongrande{Herramientas de Programación Sobre Clientes Web}

\resumen{
  La elección del editor de código depende de las necesidades del desarrollador, priorizando herramientas de desarrollo sobre editores de contenido. Es crucial utilizar validadores W3C y entender el modelo cliente-servidor para asegurar la calidad y seguridad del código web.
}

\definicion{Editor de Código}{Herramienta especializada para escribir, editar y gestionar código fuente (HTML, JavaScript, CSS), ofreciendo funcionalidades como resaltado de sintaxis y verificación de errores.}

\subtituloA{Pautas para la selección de editores}

\lista{
  **Buscar facilidad para el desarrollo web:** codificación de colores, verificación de sintaxis, inserción de comentarios y soporte FTP |
  **Ejemplos de editores gratuitos:** Notepad++, Eclipse, NetBeans, KompoZer, Bluefish.
}
